\section{Tools}
\label{sec:tools}

Everything in this section is standard. It is collected so that the later sections can refer to one place and so that the paper reads without outside references.

\subsection{Norms}

For $x\in\R^n$, $\norm x=\bigl(\sum_i x_i^2\bigr)^{1/2}$ and $\langle x,y\rangle=x^Ty$; the Cauchy--Schwarz inequality $\abs{\langle x,y\rangle}\le\norm x\norm y$ says a cosine is at most one. The unit sphere is $S^{n-1}=\{x:\norm x=1\}$. For a linear map $T:\R^n\to\R^m$ the operator norm is
\[
\op{T}=\sup_{x\ne 0}\frac{\norm{Tx}}{\norm x}=\sup_{\norm x=1}\norm{Tx},
\]
the two forms agreeing because $x=\norm x\cdot x/\norm x$. Thus $\norm{Tx}\le\op T\norm x$ for every $x$. The supremum is attained on the compact sphere; a unit vector $v$ with $\norm{Tv}=\op T$ is a right singular vector for the largest singular value. If $Q$ is orthogonal, $Q^TQ=I$, then $\norm{Qx}^2=x^TQ^TQx=\norm x^2$, so $\op{QT}=\op T=\op{TQ}$. A diagonal matrix with entries $\pm1$ is orthogonal; this is the one fact about orthogonal matrices the proof uses.

\subsection{ReLU}

$\phi(t)=\max\{t,0\}$, applied coordinatewise to vectors.

\begin{lemma}\label{lem:relu}
For all real $a,b$ and every $c\ge 0$: $\abs{\phi(a)-\phi(b)}\le\abs{a-b}$, $\phi(ca)=c\phi(a)$, and $\phi(a)=\one_{\{a>0\}}a$ whenever $a\ne 0$. Consequently $\norm{\phi(x)-\phi(y)}\le\norm{x-y}$ for vectors.
\end{lemma}

\begin{proof}
If $a,b\ge 0$ the first inequality is an equality; if $a,b\le 0$ both values are zero; if $a\ge 0\ge b$ then $\abs{\phi(a)-\phi(b)}=a\le a-b$. Homogeneity holds because multiplying by $c\ge 0$ does not change the sign of $a$, and the indicator formula is the definition away from zero. For vectors, square the scalar inequality coordinatewise and sum.
\end{proof}

Homogeneity is what makes a bias-free network positively homogeneous: from \cref{eq:forward}, $x_\ell(cs)=cx_\ell(s)$ for $c\ge 0$ by induction on $\ell$, so
\begin{equation}\label{eq:homog}
F^{(\alpha)}_{n,L}(cs)=cF^{(\alpha)}_{n,L}(s),\qquad c\ge 0,\qquad\text{and}\qquad F^{(\alpha)}_{n,L}(0)=0.
\end{equation}

\subsection{Lipschitz constants and Jacobians}

$F:\R^n\to\R^m$ is globally Lipschitz if $\norm{F(s)-F(t)}\le K\norm{s-t}$ for all $s,t$; the least such $K$ is $\Lip_2(F)$. If $F$ is differentiable at $s$ with Jacobian $DF(s)$, then for a unit vector $v$,
\[
\norm{DF(s)v}=\lim_{h\to 0}\frac{\norm{F(s+hv)-F(s)}}{\abs h}\le\Lip_2(F),
\]
so $\op{DF(s)}\le\Lip_2(F)$ at every point of differentiability. For a continuous map that is linear on each of finitely many polyhedral pieces the reverse also holds, and equality is exactly what the network needs; that is \cref{lem:pl} below.

The constant $K_{n,L}(\alpha)$ of \cref{eq:Kdef} is finite: ReLU is $1$-Lipschitz, so
\[
\norm{x_\ell(s)-x_\ell(t)}\le\op{W_\ell}\norm{x_{\ell-1}(s)-x_{\ell-1}(t)},
\]
and iterating gives $K_{n,L}(\alpha)\le\prod_\ell\op{W_\ell}<\infty$. It is measurable: by continuity the supremum in \cref{eq:Kdef} may be restricted to rational pairs, and each ratio is a Borel function of the weights.

\subsection{Three probability inequalities}

\begin{lemma}\label{lem:prob}
(a) Union bound: $\PP(\bigcup_j A_j)\le\sum_j\PP(A_j)$. (b) Markov: if $Y\ge 0$ and $a>0$ then $\PP(Y>a)\le\E Y/a$; in particular $\PP(Y>1)\le\E Y$. (c) Chebyshev: if $Y$ has finite variance then $\PP(\abs{Y-\E Y}>\varepsilon)\le\Var(Y)/\varepsilon^2$.
\end{lemma}

\begin{proof}
(a) $\one_{\cup A_j}\le\sum\one_{A_j}$ pointwise; take expectations. (b) $Y\ge a\one_{\{Y>a\}}$ pointwise; take expectations and divide by $a$. (c) Apply (b) to $(Y-\E Y)^2$ with threshold $\varepsilon^2$.
\end{proof}

\subsection{Gaussian facts}

$Z\sim\mathcal N(0,1)$ has density $(2\pi)^{-1/2}e^{-z^2/2}$, $\E Z^2=1$, $\E Z^4=3$ (two integrations by parts), and for $t<1/2$
\begin{equation}\label{eq:mgf}
\E e^{tZ^2}=\frac{1}{\sqrt{2\pi}}\int_\R e^{-(1-2t)z^2/2}\dd z=(1-2t)^{-1/2}.
\end{equation}
If $G\in\R^m$ has independent $\mathcal N(0,1)$ coordinates, $\norm G^2$ has the chi-square law with $m$ degrees of freedom, written $\chi^2_m$; we set $\chi^2_0=0$. If $Q$ is a fixed orthogonal matrix then $QG$ is again centred Gaussian with covariance $QQ^T=I$, so $QG\eqd G$; in particular multiplying rows or columns of a Gaussian matrix by fixed signs does not change its law. A centred Gaussian vector is determined by its covariance, so jointly Gaussian variables with zero covariance are independent.

\subsection{A bound on partial binomial sums}

\begin{lemma}\label{lem:entropy}
For integers $N\ge1$ and $0\le k\le N/2$,
\[
\sum_{j=0}^{k}\binom Nj\le\exp\bigl(N\,H(k/N)\bigr).
\]
\end{lemma}

\begin{proof}
For $k=0$ both sides are $1$. Otherwise put $q=k/N\in(0,1/2]$. Expanding $1=(q+(1-q))^N$ and keeping only the terms $j\le k$,
\[
1\ge\sum_{j\le k}\binom Nj q^j(1-q)^{N-j}\ge\sum_{j\le k}\binom Nj q^k(1-q)^{N-k},
\]
because $q\le 1-q$ makes $q^j(1-q)^{N-j}$ non-increasing in $j$. Hence $\sum_{j\le k}\binom Nj\le q^{-k}(1-q)^{-(N-k)}=\exp(NH(q))$.
\end{proof}

This is the only counting estimate the finite theorem needs. Stirling's formula appears in \cref{sec:sharp}, where the exact exponent is computed, and nowhere else.

\subsection{One deterministic inequality}

\begin{lemma}\label{lem:xp}
For every $x\ge 0$, $p>0$ and $t>0$, $x^p\le(p/(et))^pe^{tx}$.
\end{lemma}

\begin{proof}
For $x>0$, $\log(x^pe^{-tx})=p\log x-tx$ has derivative $p/x-t$, positive before $x=p/t$ and negative after, so the maximum of $x^pe^{-tx}$ is $(p/t)^pe^{-p}$. Multiply by $e^{tx}$.
\end{proof}

The inequality is tight at exactly one point, $x=p/t$; it is what converts a moment of order $p$ into a moment-generating function, at the price of the factor $(p/(et))^p$.
